\section{Condensate Depletion and Excitations}

We now want to compute the non condensate fraction of the weakly interacting Bose gas, which is given by the expectation value of the number operator for non-zero momentum states:
\[
    N - N_0 = \sum_{\mathbf{k}}^{\prime} \langle a_{\mathbf{k}}^\dagger a_{\mathbf{k}} \rangle,
\]
recalling that the expectation value of \(n_{\mathbf{k}} \) (number operator of the Bogoliubov quasiparticles) is \(v_{\mathbf{k}}^2\), which in the non interacting limit goes to zero:
\[
    v_{\mathbf{k}}^2 = \frac{1}{2}\left(\frac{\epsilon_k + nU_0}{E_k} - 1\right) \xrightarrow{U_0 \to 0} 0.
\]

[...]

sum to integrals

    [...]

We will just check the convergence of the integral in the ultraviolet limit (\(k \to \infty\) ), and in the infrared limit (\(k \to 0\)).

    [...]

dimensionless variables and occupation
\[
    n^{\prime} = \dots
\]
[...]

\TODO{add graph of ground state occupation number}

We can thus interact with the condensate by operating on the density, the only free parameter in \(c = \sqrt{\frac{n U_0}{m}} \).

\begin{example}[Scattering Length]
    The scattering length \(a\) is a measure of the strength of the interaction between two particles. It is defined as the distance at which the wave function of two interacting particles goes to zero. In the context of Bose-Einstein condensates, the scattering length plays a crucial role in determining the properties of the condensate, including its stability and the nature of its excitations.

    For weakly interacting Bose gases, the scattering length can be related to the interaction potential \(U_0\) through the relation:
    \[
        U_0 = \frac{4\pi \hbar^2 a}{m},
    \]
    where \(m\) is the mass of the particles. This relation allows us to express various physical quantities in terms of the scattering length, providing a more intuitive understanding of how interactions affect the condensate.

        [...]

\end{example}

condensate depletion computation with na

comment for which the result does not have any hbar term or terms we can expand to any classical or semiclassical limit. This is a consequence of the fact that the depletion and BEC in general are purely quantum effects, arising from the interactions between particles in the condensate.

    [...]

\section{Ground State Energy}

 [...]

\subsubsection{Pressure and Speed of Sound}

[...]

\subsection{Ground State Wave Function}

[...]

